% =====================================================================
%  Reporte - Proyecto Integrador: Pipeline CI/CD (Flask, Docker,
%  GitHub Actions y AWS EC2)
%  Compilar con XeLaTeX / LuaLaTeX / Tectonic (o en Overleaf con XeLaTeX)
% =====================================================================
\documentclass[12pt,letterpaper]{article}

% ---------------- DATOS DE LA PORTADA (editar aquí) ----------------
\newcommand{\institucion}{UNIVERSIDAD TECNOLÓGICA}             % <-- nombre completo de tu universidad
\newcommand{\carrera}{Ingeniería en Desarrollo y Gestión de Software}
\newcommand{\materia}{DevOps}                                  % <-- nombre oficial de la materia
\newcommand{\profesor}{Emmanuel Martínez Hernández}
\newcommand{\grupo}{IDGS16}
\newcommand{\integrantes}{Diego Garcia Martínez\\ % <-- agregar al resto del equipo
}
\newcommand{\fecha}{5 de octubre de 2026}
% Logo: colocar el archivo en figuras/logo.png
% -------------------------------------------------------------------

\usepackage[spanish,es-tabla,es-nodecimaldot]{babel}
\usepackage{fontspec}
\usepackage[margin=2.5cm]{geometry}
\usepackage{graphicx}
\usepackage{float}
\usepackage[font=small,labelfont=bf]{caption}
\usepackage{booktabs}
\usepackage{tabularx}
\usepackage{xcolor}
\usepackage{listings}
\usepackage{setspace}
\usepackage{indentfirst}
\usepackage{tikz}
\usetikzlibrary{positioning,arrows.meta,shapes.geometric,fit,backgrounds}
\usepackage[hidelinks]{hyperref}
\usepackage{url}

\onehalfspacing
\setlength{\parskip}{0.4em}
\graphicspath{{figuras/}}
\setlength{\emergencystretch}{3em}

\definecolor{codebg}{RGB}{246,246,246}
\definecolor{terminal}{RGB}{30,30,30}
\lstset{
  basicstyle=\ttfamily\footnotesize,
  backgroundcolor=\color{codebg},
  frame=single, rulecolor=\color{gray!50},
  breaklines=true, columns=fullflexible, keepspaces=true,
  showstringspaces=false, upquote=true,
  literate={á}{{\'a}}1 {é}{{\'e}}1 {í}{{\'i}}1 {ó}{{\'o}}1 {ú}{{\'u}}1 {ñ}{{\~n}}1
           {Á}{{\'A}}1 {É}{{\'E}}1 {Í}{{\'I}}1 {Ó}{{\'O}}1 {Ú}{{\'U}}1 {→}{{$\to$}}1
}
\lstdefinestyle{term}{
  backgroundcolor=\color{terminal}, basicstyle=\ttfamily\scriptsize\color{white},
  rulecolor=\color{terminal}
}

% Captura de pantalla: si el archivo existe se inserta; si no, se dibuja un
% recuadro indicando qué captura falta.  \captura{archivo}{pie de figura}{etiqueta}
\newcommand{\captura}[3]{%
  \begin{figure}[H]
    \centering
    \IfFileExists{figuras/#1}{%
      \includegraphics[width=0.92\textwidth,height=0.42\textheight,keepaspectratio]{#1}%
    }{%
      \fbox{\parbox[c][4.5cm][c]{0.85\textwidth}{\centering\color{red!70!black}
        \textbf{INSERTAR CAPTURA}\\[0.4em]\texttt{\detokenize{figuras/#1}}\\[0.4em]\small #2}}%
    }
    \caption{#2}
    \label{#3}
  \end{figure}}

\begin{document}

% =====================================================================
%  PORTADA
% =====================================================================
\begin{titlepage}
  \centering
  \IfFileExists{figuras/logo.png}{\includegraphics[height=3.2cm]{logo.png}\par\vspace{0.8cm}}{}
  {\Large\bfseries \institucion\par}
  \vspace{0.3cm}
  {\large \carrera\par}
  \vspace{2.2cm}
  {\large Materia: \materia\par}
  \vspace{1.2cm}
  \rule{\textwidth}{1pt}\par\vspace{0.4cm}
  {\LARGE\bfseries Proyecto Integrador:\par}
  \vspace{0.3cm}
  {\Large\bfseries Pipeline CI/CD automatizado para una API REST con Docker, GitHub Actions y AWS EC2\par}
  \vspace{0.4cm}\rule{\textwidth}{1pt}\par
  \vspace{2cm}
  {\large \textbf{Integrantes:}\par}
  \vspace{0.2cm}
  {\large \integrantes\par}
  \vspace{1cm}
  {\large \textbf{Grupo:} \grupo\par}
  \vspace{0.3cm}
  {\large \textbf{Profesor:} \profesor\par}
  \vfill
  {\large \fecha\par}
\end{titlepage}

\tableofcontents
\newpage
\listoffigures
\newpage

% =====================================================================
\section{Introducción}
% =====================================================================

El desarrollo de software moderno ya no termina cuando el programador escribe la última línea de
código. Para que una aplicación genere valor tiene que llegar a producción de forma rápida,
repetible y segura, y tiene que poder actualizarse muchas veces sin que los usuarios lo noten.
Durante años el paso del código al servidor se hacía de forma manual: alguien compilaba el
proyecto en su computadora, copiaba archivos al servidor, reiniciaba servicios y cruzaba los dedos.
Este proceso era lento, dependía de la memoria de una persona y era una fuente constante de
errores del tipo ``en mi máquina sí funciona''. La cultura DevOps surge precisamente para cerrar
la brecha entre quienes desarrollan y quienes operan el software, apoyándose en la automatización
de todo el ciclo de vida de una aplicación \cite{kim2016}.

Dos prácticas son el corazón de esa automatización. La \textbf{integración continua} (CI,
\textit{Continuous Integration}) consiste en que cada cambio que se sube al repositorio se
compila y se prueba automáticamente, de modo que los errores se detectan minutos después de
haberse introducido y no semanas después \cite{fowler2006}. La \textbf{entrega y despliegue
continuo} (CD, \textit{Continuous Delivery/Deployment}) extiende esa idea: si las pruebas pasan,
el mismo proceso empaqueta la aplicación y la publica en el servidor sin intervención humana
\cite{humble2010}. Al conjunto ordenado de pasos que sigue el código desde el \texttt{git push}
hasta producción se le llama \textit{pipeline}.

Para que el pipeline sea confiable, el artefacto que se prueba debe ser exactamente el mismo que
se ejecuta en producción. Aquí entran los \textbf{contenedores}: Docker permite empaquetar la
aplicación junto con su intérprete y sus dependencias en una imagen inmutable que se comporta igual
en la laptop del desarrollador, en el servidor de integración y en la nube \cite{dockerdocs}.
Estas imágenes se guardan en un registro (en este proyecto, Docker Hub) y se identifican con
etiquetas; usar el hash del commit como etiqueta permite saber con precisión qué versión del código
está corriendo y regresar a una versión anterior si algo falla.

En este proyecto integrador se implementó un flujo de trabajo completo y profesional para el ciclo
de vida de una API REST. Se desarrolló una API en Python con el microframework Flask
\cite{flaskdocs} que expone 6 endpoints para la gestión de usuarios: las cinco operaciones
CRUD (listar, consultar, crear, actualizar y eliminar) y un endpoint de salud que informa la
versión desplegada. La API cuenta con 31 pruebas automatizadas escritas con pytest, que alcanzan una
cobertura de código del 99\,\%, muy por encima del umbral mínimo exigido del 70\,\%. El pipeline,
definido con GitHub Actions \cite{githubactions}, ejecuta el análisis estático y las pruebas en cada
\textit{push} o \textit{pull request} a la rama \texttt{main}; cuando el cambio llega a
\texttt{main}, construye la imagen Docker, la publica en Docker Hub con las etiquetas
\texttt{:latest} y \texttt{:<sha del commit>}, y finalmente se conecta por SSH a una instancia
Ubuntu Server de Amazon EC2 \cite{awsec2} para descargar la versión nueva, validarla, detener el
contenedor antiguo y levantar el nuevo en el puerto 80.

Un aspecto central del proyecto es la seguridad: ninguna contraseña, token, dirección IP ni llave
SSH aparece en el código del repositorio. Todos estos datos se administran como
\textit{GitHub Secrets}, y la autenticación con Docker Hub se realiza mediante un
\textit{Personal Access Token} en lugar de la contraseña de la cuenta. Además, el despliegue
incluye una verificación de salud: la versión nueva se prueba en un puerto temporal antes de
reemplazar a la que está en producción, de forma que una imagen defectuosa nunca deja al servicio
fuera de línea.

El presente reporte documenta la arquitectura de la solución, el proceso seguido en cada etapa y
los resultados obtenidos, apoyados con figuras numeradas que muestran la evidencia de cada paso,
incluyendo la demostración en vivo en la que un cambio de una línea en la API recorre
automáticamente todo el pipeline hasta quedar publicado en la nube.

% =====================================================================
\section{Objetivos}
% =====================================================================

\subsection{Objetivo general}
Implementar un pipeline CI/CD automatizado que pruebe, empaquete y despliegue una API REST en
AWS EC2 cada vez que se publica un cambio en la rama \texttt{main} del repositorio.

\subsection{Objetivos específicos}
\begin{itemize}
  \item Desarrollar una API REST con al menos 6 endpoints funcionales.
  \item Implementar pruebas unitarias y de integración con una cobertura mínima del 70\,\%.
  \item Contenedorizar la aplicación con un \texttt{Dockerfile} optimizado y un \texttt{.dockerignore}.
  \item Automatizar pruebas, construcción y publicación de la imagen con GitHub Actions y Docker Hub.
  \item Desplegar automáticamente en una instancia EC2 por SSH, gestionando los datos sensibles con GitHub Secrets.
\end{itemize}

\subsection{Entregables}
\begin{itemize}
  \item \textbf{Repositorio de GitHub:} \url{https://github.com/2023371031-del/sa-cicd-api}
  \item \textbf{URL pública de la API:} \url{http://18.223.115.56/api/health}
\end{itemize}

% =====================================================================
\section{Herramientas utilizadas}
% =====================================================================

\begin{table}[H]
\centering
\caption{Tecnologías utilizadas en el proyecto}
\label{tab:herramientas}
\begin{tabularx}{\textwidth}{@{}lX@{}}
\toprule
\textbf{Herramienta} & \textbf{Uso en el proyecto} \\
\midrule
Python 3.12 + Flask 3.1 & Lenguaje y framework de la API REST \\
Gunicorn 23 & Servidor WSGI de producción dentro del contenedor \\
pytest + pytest-cov & Pruebas automatizadas y reporte de cobertura \\
flake8 & Análisis estático (lint) del código \\
Docker & Construcción de la imagen y ejecución del contenedor \\
Docker Hub & Registro de imágenes (artefactos) \\
GitHub / GitHub Actions & Repositorio de código y motor del pipeline CI/CD \\
AWS EC2 (Ubuntu Server 26.04, \texttt{t3.micro}) & Servidor en la nube donde se ejecuta la API \\
Visual Studio Code, curl & Edición de código y pruebas manuales de endpoints \\
\bottomrule
\end{tabularx}
\end{table}

% =====================================================================
\section{Arquitectura de la solución}
% =====================================================================

La Figura~\ref{fig:arquitectura} resume el flujo completo. El desarrollador hace \texttt{git push}
a la rama \texttt{main}; GitHub Actions ejecuta tres \textit{jobs} encadenados (\texttt{test},
\texttt{docker} y \texttt{deploy}), cada uno de los cuales solo inicia si el anterior terminó con
éxito. El resultado final es un contenedor actualizado en la instancia EC2 que atiende peticiones
HTTP en el puerto 80.

\begin{figure}[H]
\centering
\begin{tikzpicture}[
  node distance=0.7cm and 0.6cm,
  box/.style={draw, rounded corners, align=center, minimum height=1.1cm, text width=2.6cm, font=\footnotesize, fill=white},
  job/.style={box, fill=blue!8},
  ext/.style={box, fill=orange!12},
  aws/.style={box, fill=green!10},
  arr/.style={-{Stealth[length=2.2mm]}, thick}
]
  \node[box] (dev) {Desarrollador\\\texttt{git push}};
  \node[ext, right=of dev] (gh) {GitHub\\rama \texttt{main}};
  \node[job, right=of gh] (test) {\textbf{job: test}\\flake8 + pytest\\cobertura $\geq$ 70\,\%};
  \node[job, below=1.1cm of test] (docker) {\textbf{job: docker}\\build + push\\\texttt{:latest}, \texttt{:sha}};
  \node[ext, left=of docker] (hub) {Docker Hub\\\texttt{sa-cicd-api}};
  \node[job, below=1.1cm of docker] (deploy) {\textbf{job: deploy}\\SSH (.pem)\\pull + health check};
  \node[aws, left=of deploy] (ec2) {AWS EC2 Ubuntu\\contenedor \texttt{api}\\puerto 80};
  \node[box, left=of ec2] (user) {Usuario\\{\scriptsize\texttt{http://IP/api/...}}};
  \node[ext, right=0.5cm of docker, text width=2.1cm] (sec) {GitHub\\Secrets};

  \begin{scope}[on background layer]
    \node[draw, dashed, rounded corners, fit=(test)(docker)(deploy), inner sep=6pt,
          label={[font=\footnotesize]above:GitHub Actions}] {};
  \end{scope}

  \draw[arr] (dev) -- (gh);
  \draw[arr] (gh) -- (test);
  \draw[arr] (test) -- (docker);
  \draw[arr] (docker) -- (hub);
  \draw[arr] (docker) -- (deploy);
  \draw[arr] (deploy) -- (ec2);
  \draw[arr, dashed] (hub) -- node[font=\scriptsize, left]{\texttt{docker pull}} (ec2);
  \draw[arr] (user) -- (ec2);
  \draw[arr, dotted] (sec) -- (docker);
  \draw[arr, dotted] (sec) |- (deploy);
\end{tikzpicture}
\caption{Arquitectura del pipeline CI/CD implementado}
\label{fig:arquitectura}
\end{figure}

El código se organizó por capas, siguiendo la estructura vista en clase: \texttt{routes/} contiene
las rutas (\texttt{users.py} y \texttt{system.py}), \texttt{services/} la lógica de negocio y las validaciones,
\texttt{data/} el almacenamiento, y \texttt{tests/} las pruebas. En la raíz se encuentran el
\texttt{Dockerfile}, el \texttt{.dockerignore} y el workflow en \texttt{.github/workflows/main.yml}.

% =====================================================================
\section{Resultados}
% =====================================================================

\subsection{Desarrollo de la API REST}

Se construyó una API de gestión de usuarios con \textbf{6 endpoints}: cinco operaciones CRUD sobre
el recurso \texttt{users} y un endpoint de salud (\texttt{/api/health}) que además informa la
versión, un mensaje configurable y el hash del commit desplegado. El código se separó en capas:
las rutas (\texttt{routes/users.py}) solo reciben la petición y devuelven la respuesta; la lógica
de negocio (\texttt{services/user\_service.py}) valida los datos contra un esquema de campos,
revisa que el email no se repita y modifica los datos; y \texttt{data/store.py} mantiene los
usuarios en memoria a partir de \texttt{data/seed.json}. La Figura~\ref{fig:estructura} muestra la
organización del proyecto y la Tabla~\ref{tab:endpoints} resume los endpoints.

\captura{estructura.png}{Estructura de archivos del repositorio}{fig:estructura}

\begin{table}[H]
\centering
\caption{Endpoints de la API}
\label{tab:endpoints}
\footnotesize
\begin{tabularx}{\textwidth}{@{}clXl@{}}
\toprule
\textbf{\#} & \textbf{Método} & \textbf{Ruta y descripción} & \textbf{Códigos} \\
\midrule
1 & GET    & \texttt{/api/health} --- estado, versión, mensaje y commit & 200 \\
2 & GET    & \texttt{/api/users} --- listar usuarios & 200 \\
3 & GET    & \texttt{/api/users/\{id\}} --- consultar un usuario & 200, 404 \\
4 & POST   & \texttt{/api/users} --- crear un usuario & 201, 400, 409, 415 \\
5 & PUT    & \texttt{/api/users/\{id\}} --- actualizar uno o varios campos & 200, 400, 404, 409, 415 \\
6 & DELETE & \texttt{/api/users/\{id\}} --- eliminar un usuario & 200, 404 \\
\bottomrule
\end{tabularx}
\end{table}

Todas las respuestas comparten el formato \texttt{\{"statusCode": ..., "data": ...\}} y los errores
\texttt{\{"statusCode": ..., "error": "mensaje"\}}, con códigos HTTP coherentes: 201 al crear, 400
ante datos inválidos (campos faltantes, vacíos, de tipo incorrecto, email mal formado o campos no
permitidos como \texttt{id}), 404 si el usuario no existe, 409 si el email ya está registrado y 415
cuando el cuerpo no es JSON. La Figura~\ref{fig:endpoints} muestra la prueba de los seis endpoints
con \texttt{curl} directamente contra la IP pública de la EC2: se consulta el estado, se listan los
usuarios, se consulta uno, se crea el usuario 4, se actualiza su rol y finalmente se elimina.

\captura{endpoints.png}{Prueba de los 6 endpoints contra la IP pública de la EC2}{fig:endpoints}

\subsection{Pruebas automatizadas y cobertura}

Las pruebas se escribieron con pytest usando el cliente de pruebas de Flask, que permite hacer
peticiones HTTP a la aplicación sin levantar un servidor. Un \textit{fixture} recarga los datos
iniciales antes de cada prueba para que sean independientes entre sí. Se implementaron 31 pruebas
en dos archivos: \texttt{test\_users.py} prueba los cinco endpoints de usuarios, tanto los casos
exitosos como los errores que puede provocar quien consume la API (email repetido o inválido,
campos vacíos, tipos incorrectos, campos no permitidos, cuerpo sin JSON, ids inexistentes y
método no permitido); para los casos de validación se usó \textit{parametrización}, de modo que una
sola función de prueba se ejecuta con nueve cuerpos inválidos distintos. Por su parte,
\texttt{test\_system.py} cubre el endpoint de salud y el manejo de errores 404 y 500. La
Figura~\ref{fig:pytest} muestra la ejecución.

\begin{figure}[H]
\centering
\begin{minipage}{0.95\textwidth}
\begin{lstlisting}[style=term]
$ pytest -v
tests/test_system.py::test_health PASSED                                 [  3%]
tests/test_system.py::test_route_not_found PASSED                        [  6%]
tests/test_system.py::test_internal_error_handler PASSED                 [  9%]
tests/test_users.py::test_list_users PASSED                              [ 12%]
tests/test_users.py::test_get_user PASSED                                [ 16%]
tests/test_users.py::test_get_user_not_found PASSED                      [ 19%]
...
tests/test_users.py::test_delete_user_not_found PASSED                   [ 96%]
tests/test_users.py::test_delete_all_not_allowed PASSED                  [100%]
============================= 31 passed in 0.57s ==============================
\end{lstlisting}
\end{minipage}
\caption{Ejecución local de las 31 pruebas con pytest}
\label{fig:pytest}
\end{figure}

La cobertura se mide con pytest-cov. En \texttt{pytest.ini} se configuró la opción
\texttt{-{}-cov-fail-under=70}, de modo que si la cobertura baja del 70\,\% el comando termina con
error y el pipeline se detiene. Como se observa en la Figura~\ref{fig:cobertura}, se obtuvo una
cobertura total del \textbf{99.48\,\%}: de 191 sentencias solo una no se ejecuta.

\begin{figure}[H]
\centering
\begin{minipage}{0.95\textwidth}
\begin{lstlisting}[style=term]
Name                       Stmts   Miss  Cover   Missing
--------------------------------------------------------
app.py                        24      0   100%
config.py                      2      0   100%
data/__init__.py               0      0   100%
data/store.py                 25      0   100%
routes/__init__.py             3      0   100%
routes/responses.py           13      0   100%
routes/system.py               8      0   100%
routes/users.py               20      0   100%
services/__init__.py           0      0   100%
services/user_service.py      46      0   100%
services/validators.py        50      1    98%   33
--------------------------------------------------------
TOTAL                        191      1    99%
Required test coverage of 70% reached. Total coverage: 99.48%
31 passed in 0.57s
\end{lstlisting}
\end{minipage}
\caption{Reporte de cobertura de código (umbral mínimo: 70\,\%)}
\label{fig:cobertura}
\end{figure}

\subsection{Contenedorización}

El \texttt{Dockerfile} (Figura~\ref{fig:dockerfile}) usa una construcción en dos etapas
(\textit{multi-stage build}): la primera instala las dependencias en un entorno virtual y la
segunda copia únicamente ese entorno y el código a una imagen \texttt{python:3.12-slim} limpia,
sin compiladores ni cachés de pip. Además, el contenedor se ejecuta con un usuario sin privilegios
(\texttt{appuser}), define un \texttt{HEALTHCHECK} contra \texttt{/api/health} y recibe el hash del
commit como argumento de construcción para exponerlo en \texttt{/api/health}. Las dependencias se
copian antes que el código para aprovechar la caché de capas: si solo cambia el código, no se
reinstalan las librerías \cite{dockerdocs}.

\begin{figure}[H]
\centering
\begin{minipage}{0.95\textwidth}
\begin{lstlisting}
FROM python:3.12-slim AS builder
RUN python -m venv /opt/venv
ENV PATH="/opt/venv/bin:$PATH"
COPY requirements.txt .
RUN pip install --no-cache-dir -r requirements.txt

FROM python:3.12-slim
ENV PYTHONDONTWRITEBYTECODE=1 PYTHONUNBUFFERED=1 PATH="/opt/venv/bin:$PATH"
RUN useradd --create-home --uid 1000 appuser
WORKDIR /app
COPY --from=builder /opt/venv /opt/venv
COPY --chown=appuser:appuser . .
ARG GIT_SHA=local
ENV GIT_SHA=$GIT_SHA
USER appuser
EXPOSE 8000
HEALTHCHECK --interval=30s --timeout=3s CMD python -c "import urllib.request; \
  urllib.request.urlopen('http://127.0.0.1:8000/api/health', timeout=2)" || exit 1
CMD ["gunicorn", "--bind", "0.0.0.0:8000", "--workers", "1", "--threads", "4", "app:create_app()"]
\end{lstlisting}
\end{minipage}
\caption{Dockerfile multi-etapa de la API}
\label{fig:dockerfile}
\end{figure}

El archivo \texttt{.dockerignore} excluye de la imagen todo lo que no se necesita en producción:
el directorio \texttt{.git}, cachés de Python, entornos virtuales, las pruebas, reportes de
cobertura, archivos \texttt{.env}, llaves \texttt{.pem}, logs, la documentación y el reporte.
Gracias a esto, la imagen final ocupa 49.4\,MB comprimida. La Figura~\ref{fig:dockerlocal} muestra
la construcción y prueba local del contenedor.

\begin{figure}[H]
\centering
\begin{minipage}{0.95\textwidth}
\begin{lstlisting}[style=term]
$ docker build -t sa-cicd-api --build-arg GIT_SHA=test123 .
 => naming to docker.io/library/sa-cicd-api:latest                    done
$ docker run -d --name sa-test -p 8089:8000 sa-cicd-api
$ curl http://localhost:8089/api/health
{"statusCode":200,"data":{"status":"ok","version":"1.0.0",
 "message":"Hola desde el pipeline CI/CD - IDGS16","commit":"test123"}}
$ docker images sa-cicd-api
IMAGE                ID             DISK USAGE   CONTENT SIZE
sa-cicd-api:latest   639a0ff38f75        205MB         49.4MB
\end{lstlisting}
\end{minipage}
\caption{Construcción y ejecución local de la imagen Docker}
\label{fig:dockerlocal}
\end{figure}

\subsection{Integración continua con GitHub Actions}

El workflow \texttt{.github/workflows/main.yml} se dispara con eventos \texttt{push} y
\texttt{pull\_request} sobre la rama \texttt{main}. El primer job, \texttt{test}, instala
Python 3.12, ejecuta \texttt{flake8} y luego \texttt{pytest -v}, que imprime en el log la tabla de
cobertura; adicionalmente, el resumen de cobertura se publica en formato Markdown en la página de
la ejecución (\texttt{\$GITHUB\_STEP\_SUMMARY}) y el archivo \texttt{coverage.xml} se guarda como
artefacto. En los \textit{pull requests} solo corre este job; los de construcción y despliegue se
ejecutan únicamente con \texttt{push} a \texttt{main}. La Figura~\ref{fig:actions} muestra una
ejecución exitosa del pipeline y la Figura~\ref{fig:actionscov} el resumen de cobertura en los
logs.

\captura{actions_pipeline.png}{Ejecución exitosa del pipeline con los jobs test, docker y deploy}{fig:actions}

\captura{actions_cobertura.png}{Resumen de cobertura mostrado en los logs de GitHub Actions}{fig:actionscov}

\subsection{Gestión de artefactos en Docker Hub}

El job \texttt{docker} se autentica en Docker Hub con la acción \texttt{docker/login-action}
usando un \textit{Personal Access Token} con permisos de lectura y escritura, almacenado en el
secreto \texttt{DOCKERHUB\_TOKEN}. Después construye la imagen con Buildx (con caché de capas entre
ejecuciones) y la publica con dos etiquetas: \texttt{:latest}, que siempre apunta a la versión más
reciente, y \texttt{:\$\{\{ github.sha \}\}}, que identifica de forma única el commit que la generó
y permite regresar a cualquier versión anterior. La Figura~\ref{fig:dockerhub} muestra el
repositorio en Docker Hub con ambas etiquetas.

\captura{dockerhub_tags.png}{Repositorio en Docker Hub con las etiquetas \texttt{latest} y el hash del commit}{fig:dockerhub}

\subsection{Infraestructura en AWS EC2}

Se utilizó una instancia EC2 \texttt{t3.micro} con Ubuntu Server 26.04 LTS en la región
\texttt{us-east-2} (Ohio), con IP pública \texttt{18.223.115.56}, a la que se accede con un par de
llaves en formato \texttt{.pem} (Figura~\ref{fig:ec2}). Se reutilizó la instancia de la práctica
anterior para no consumir créditos adicionales; ya contaba con Docker Engine instalado y con el
usuario \texttt{ubuntu} dentro del grupo \texttt{docker}. Para preparar una instancia nueva, el
repositorio incluye el script \texttt{scripts/setup\_ec2.sh}, que instala Docker y configura ese
grupo. Al \textit{Security Group} se le agregó la regla de entrada HTTP en el puerto 80, por donde se
atiende a los usuarios, y se abrió el puerto 22 (SSH) a cualquier origen, ya que GitHub Actions se
conecta desde direcciones IP que cambian en cada ejecución (Figura~\ref{fig:sg}).

\captura{ec2_instancia.png}{Instancia EC2 en ejecución en la consola de AWS}{fig:ec2}

\captura{ec2_security_group.png}{Reglas de entrada del Security Group (puertos 22 y 80)}{fig:sg}

\subsection{Seguridad: GitHub Secrets}

Ningún dato sensible se escribió en el código. En el repositorio se configuraron cinco secretos
(Figura~\ref{fig:secrets}): \texttt{DOCKERHUB\_USERNAME}, \texttt{DOCKERHUB\_TOKEN},
\texttt{EC2\_HOST} (IP pública de la instancia), \texttt{EC2\_USER} y \texttt{EC2\_SSH\_KEY}
(contenido del archivo \texttt{.pem}). GitHub los cifra, los inyecta solo durante la ejecución y
los oculta (\texttt{***}) en los logs. Como medida adicional, \texttt{.gitignore} y
\texttt{.dockerignore} excluyen archivos \texttt{.env}, \texttt{*.pem} y \texttt{*.key}.

\captura{github_secrets.png}{Secretos configurados en el repositorio (GitHub nunca muestra sus valores)}{fig:secrets}

\subsection{Despliegue continuo en EC2}

El job \texttt{deploy} se conecta a la instancia por SSH con \texttt{appleboy/ssh-action} y
ejecuta el siguiente proceso:

\begin{enumerate}
  \item Descarga la imagen \texttt{:latest} recién publicada en Docker Hub.
  \item Levanta la versión nueva en un contenedor temporal (puerto 8081) y consulta
        \texttt{/api/health} hasta 20 veces. Si no responde, muestra sus logs, lo elimina y
        marca el pipeline como fallido \textbf{sin tocar} la versión en producción.
  \item Si está sana, elimina el contenedor temporal, detiene el contenedor antiguo
        (\texttt{api}) y levanta la nueva versión mapeada al puerto 80 con
        \texttt{--restart unless-stopped}.
  \item Verifica que \texttt{/api/health} responda con el hash del commit del \textit{push}, lo
        que comprueba que realmente se está sirviendo la versión nueva, y limpia imágenes viejas.
\end{enumerate}

Como la imagen ya está descargada y validada antes de detener la versión anterior, el cambio de
contenedor tarda alrededor de un segundo, reduciendo el tiempo de inactividad al mínimo. En las
pruebas realizadas, el job completo de despliegue tardó unos 10 segundos. La
Figura~\ref{fig:deploylog} muestra el log del despliegue, la Figura~\ref{fig:dockerps} el
contenedor corriendo en la EC2 y la Figura~\ref{fig:apipublica} la API respondiendo en la IP
pública.

\captura{actions_deploy_log.png}{Log del job deploy en GitHub Actions}{fig:deploylog}

\captura{ec2_docker_ps.png}{Contenedor \texttt{api} en ejecución en la EC2 (\texttt{docker ps})}{fig:dockerps}

\captura{api_publica.png}{API respondiendo en el navegador en \texttt{http://18.223.115.56/api/health}}{fig:apipublica}

\subsection{Demostración práctica}

Para la demostración en vivo se modificó una sola línea del archivo \texttt{config.py}, la
constante \texttt{MESSAGE}, y se hizo \texttt{git commit} y \texttt{git push} a \texttt{main}. Antes
del cambio, \texttt{GET /api/health} mostraba el mensaje ``Hola desde el pipeline CI/CD - IDGS16''
y el commit \texttt{47737e5} (Figura~\ref{fig:antes}). El \textit{push} disparó la ejecución
\#6 del pipeline (Figura~\ref{fig:actions}), que en 59 segundos ejecutó las pruebas, publicó la
etiqueta \texttt{1a63431...} en Docker Hub (Figura~\ref{fig:dockerhub}) y actualizó el servidor. Al
volver a consultar la URL se obtuvo el mensaje nuevo, ``Demo en vivo: versión 2 desplegada
automáticamente - IDGS16'', junto con el commit \texttt{1a63431} (Figura~\ref{fig:despues}), sin
haber entrado manualmente al servidor.

\captura{demo_antes.png}{\texttt{/api/health} antes del cambio, y el \texttt{git commit} / \texttt{git push} de la demo}{fig:antes}

\captura{demo_despues.png}{\texttt{/api/health} después del \texttt{git push}: mensaje y commit nuevos}{fig:despues}

% =====================================================================
\section{Conclusiones}
% =====================================================================

El proyecto permitió construir de principio a fin un flujo de trabajo CI/CD como el que se usa en
la industria. Se desarrolló una API REST con 6 endpoints, organizada en capas de rutas, servicios
y datos, y respaldada por 31 pruebas automatizadas con una cobertura del 99.48\,\%. Más allá del
número, las pruebas resultaron ser la pieza que da confianza a todo el pipeline: como el umbral de
cobertura y el resultado de las pruebas son condiciones obligatorias para continuar, ningún cambio
que rompa la API puede llegar a Docker Hub ni al servidor. Separar el código en rutas, servicios
y datos también demostró su valor: las reglas de negocio (validaciones, email único) viven en un
solo lugar, lo que facilita probarlas a fondo y agregar nuevos recursos en el futuro sin tocar la
lógica existente.

La contenedorización con Docker resolvió el problema de las diferencias entre entornos: la misma
imagen que se probó localmente es la que se publica y se ejecuta en EC2, identificada con el hash
de su commit, lo que facilita saber qué versión está en producción y regresar a una anterior si es
necesario. GitHub Actions unió todas las piezas en un solo proceso reproducible y visible para
todo el equipo, y GitHub Secrets permitió automatizar el acceso a Docker Hub y a AWS sin exponer
contraseñas, tokens, IPs ni llaves en un repositorio público, que es uno de los errores de
seguridad más comunes en proyectos reales.

Finalmente, la demostración práctica evidenció el beneficio principal de DevOps: un cambio de una
línea pasó de la computadora del desarrollador a producción en pocos minutos, sin intervención
manual y con verificaciones automáticas de salud antes y después del despliegue. Como trabajo
futuro, la solución podría incorporar una base de datos persistente (por ejemplo, Amazon RDS),
un proxy inverso con despliegue \textit{blue-green} para eliminar por completo el segundo de
inactividad durante el cambio de contenedor, y HTTPS mediante un certificado TLS.

% =====================================================================
\renewcommand{\refname}{Fuentes de información}
\begin{thebibliography}{9}
% =====================================================================

\bibitem{dockerdocs}
Docker Inc. (2026). \textit{Building best practices}. Docker Documentation.
\url{https://docs.docker.com/build/building/best-practices/}

\bibitem{flaskdocs}
Pallets Projects. (2026). \textit{Flask documentation (3.1.x)}.
\url{https://flask.palletsprojects.com/}

\bibitem{fowler2006}
Fowler, M. (2006). \textit{Continuous integration}. martinfowler.com.
\url{https://martinfowler.com/articles/continuousIntegration.html}

\bibitem{githubactions}
GitHub. (2026). \textit{GitHub Actions documentation}. GitHub Docs.
\url{https://docs.github.com/en/actions}

\bibitem{humble2010}
Humble, J., \& Farley, D. (2010). \textit{Continuous delivery: Reliable software releases through
build, test, and deployment automation}. Addison-Wesley.

\bibitem{kim2016}
Kim, G., Humble, J., Debois, P., \& Willis, J. (2016). \textit{The DevOps handbook: How to create
world-class agility, reliability, and security in technology organizations}. IT Revolution Press.

\bibitem{awsec2}
Amazon Web Services. (2026). \textit{Amazon EC2 user guide}. AWS Documentation.
\url{https://docs.aws.amazon.com/ec2/}

\end{thebibliography}

\end{document}
